%% =====================================================================
%%  第 11 章　冲突
%%  源文件：11-冲突.md
%% =====================================================================
\chapter{冲突}

\applicable{你和一个人正在吵，而且越吵越大；或者你们已经不吵了，各自沉默着。}

\begin{guideblock}
\textbf{本章解决什么问题}：为什么越讲道理，对方越激动；一句“你先说”，为什么能停下来；以及吵完之后，那三件必须按顺序做的事。
\end{guideblock}

先给结论：\textbf{吵架是一次死锁。}双方同时在发送，没有人在接收。两边都在等对方先收，而先收的那个人看起来像认输。

值得注意的是，绝大多数争吵不是在争论事实，是在争夺“被听见”。所以在这个场合里，\textbf{讲道理是无效的——道理解决的是内容，而他们抢的是状态。}

\hcirule

\section{一次吵架的拆解}

先看样本。

\sample{1573}\par
\begin{mdquote}
你：「你昨天为什么不说？」\par
他：「我说了啊，你自己没听。」\par
你：「我什么时候没听——」\par
他：「你每次都这样，一有事就先怪别人。」\par
你：「我怪你了？我什么时候怪你了？」\par
他：「你看，你现在就在怪。」\par
\vspace{0.3\baselineskip}
二十分钟后，你们在吵三年前的事。
\end{mdquote}

把它逐句拆开：

\begin{manstable}{P{2.2cm} L P{2.6cm}}
\tbh{说话人} & \tbh{这句话是什么} & \tbh{是不是在接收} \\
\midrule
你 & 发送（质问） & 否 \\
他 & 发送（反质问） & 否 \\
你 & 发送（否认） & 否 \\
他 & 发送（扩大指控） & 否 \\
你 & 发送（要求举证） & 否 \\
他 & 发送（用当前行为举证） & 否 \\
\bottomrule
\end{manstable}

\textbf{六句话，零次接收。}

这就是死锁。两个人都在发，都在等对方收。越往后，发送的强度越高、内容越远，因为\textbf{音量是被“对方没有接收”这件事本身推高的}，而不是被原始问题推高的。

二十分钟后开始翻三年前的旧账，也是同一个原因：\textbf{当前这件事已经被证明无法传达}，于是系统开始寻找别的载荷，看看哪一条能送进去。

需要说明的是，这个过程不携带任何恶意。两个人都在真诚地试图把某样东西送出去，只是谁也没有做接收这件事。

再看一个升级更快的版本。这一次的主角是一对在一起三年的情侣，起手只有一句话。

\sample{1600}\par
\begin{mdquote}
他：「周末你到底去不去我妈那边？」\par
她：「我周五跟你说过这周要加班。」\par
他：「你跟我说的是‘可能’。」\par
她：「我说‘可能’，是因为我还没确定，不是给你留个话口以后来怪我。」\par
他：「我什么时候怪你了，我只是问你去不去。」\par
她：「你刚才那句‘你跟我说的是可能’——你自己听听那语气。」\par
他：「什么语气？我就正常说话。」\par
她：「行，那我也不用去了。」
\end{mdquote}

从第四句开始，争的东西已经不是“周末去哪”，是“你那句话到底什么意思”。

\textbf{这场争吵在两句话之内就完成了换题。}原始议题（周末的行程）在此之后再也没有被提起——它不是被解决了，是被挤出去了。

值得注意的是，越往后，说得越少，声音越大。\textbf{争吵末期的句子普遍比开头短三到四成。}内容被压缩，情绪被放大，这是死锁进入后期的典型特征。

\section{吵架到底在吵什么}

\hciTag{核心结论}：判断一次争吵的类型，只看一个问题——争的是“怎么办”，还是“你刚才那句话什么意思”。

\begin{manstable}{P{1.8cm} L L}
\tbh{} & \tbh{内容之争} & \tbh{状态之争} \\
\midrule
争什么 & 这件事怎么办 & 你有没有在听我 \\
判据 & 焦点在事情本身 & 焦点在对方的一句话、一个语气 \\
讲道理 & 有效 & \textbf{无效} \\
结束条件 & 有了方案 & 有一方感到被听见 \\
在样本库中的占比 & 约两成 & 约八成 \\
\bottomrule
\end{manstable}

需要说明的是，\textbf{这两类争吵的外表是一样的。}它们都用同样的话开始：“你怎么能这么说。”区别在于内容之争里，“这么说”是一个需要讨论的观点；状态之争里，“这么说”是一个需要被承认的感受。

\textbf{把状态之争当成内容之争来处理，是吵架升级最快的一条路。}你越有理，对方越激动——因为你的每一个道理都在说“你说的内容不成立”，而他要的从来不是内容成立。

\textbf{换题的瞬间，就是类型切换的信号。}当话题从“这件事”滑到“你那句话的语气”，说明它已经从内容之争切进了状态之争；此后所有关于事情的论据，都是在给一个不再开着的入口送货。

\section{死锁怎么解}

死锁的定义是：双方都在等对方先收。

\textbf{所以解法只有一条：有一个人先收。}

具体动作：

\begin{mdquote}
不反驳，让对方说完。\par
「你先说。」\par
「我听你说完。」
\end{mdquote}

这三句在内容上什么都没让步——你没有承认任何指控，也没有撤回任何观点。\textbf{它让出的只有一样东西：发送权。}

而发送权正是死锁里唯一稀缺的资源。

\textbf{它的成本是真实的：先收的人看起来像输。}这是要付的代价，本手册不打算美化它。但需要说明的是，\textbf{在死锁里没有“赢”这个结果}——双方继续发下去，得到的只是两个人一起损失时间、损失关系，以及额外增加的一批旧账。

有读者会问：如果我先收，对方趁机把话全说成他的功劳，怎么办？

答案是：那说明你要处理的已经不是这一场冲突，而是这段关系本身。\textbf{先收这个动作，同时也是在测试。“我说‘你先说’之后，对方是接着说，还是趁机定罪。”前一种关系值得继续，后一种不值得。}测试的成本，是这一场冲突里你让出的那几分钟。

\section{课题分离}

这一节的用处，是在冲突发生之前。

人类有一类冲突，是\textbf{请求被路由到了错误的处理者身上}。

\begin{mdquote}
你在公司被上级当众否定，回家之后因为伴侣没洗碗而爆发。
\end{mdquote}

拆开看：你带着一个未处理的请求（“有人不尊重我”）回到家，而这个请求\textbf{无法在伴侣身上被解决}。伴侣不是那个欠你回应的人。他接到的是一次错投。

\textbf{课题分离，就是把请求路由给正确的处理者。}

动手问自己一句：

\begin{mdquote}
\textbf{这场冲突真正的对象是谁？是眼前这个人，还是别人、别的事？}
\end{mdquote}

如果是后者，正确的做法是\textbf{不对眼前的人发起这场冲突}——你可以告诉他你遇到了什么事，但不要让他替别人承担后果。

需要强调的是，课题分离不是叫你闭口。它是叫你\textbf{分清哪一段是要处理的，哪一段是要讲的}。

\section{冲突之后}

冲突结束（或者暂停）之后，有一个高危窗口。

\textbf{不要立刻修复。}

立刻道歉、立刻解释、立刻要求对方表态，都属于抢跑——原因\hciSee{第12章}。此时双方的情绪值都还在高位，任何修复动作都会被解读成“想尽快结束这个话题”，而不是“想解决这个问题”。

但另一端也要防：\textbf{不要用长时间沉默当武器。}冷处理的边界是——它是为了让情绪衰减，不是为了惩罚对方。前者有明确的长度，后者没有。

需要提醒的是，冷处理有一个\textbf{长度上限}：通常在一天到两天之间。超过之后，它不再是“让情绪衰减”，而是变成了 11.10 要讲的冷战。

\section{冷战的运行机制}

吵到某个点，双方会同时停下来。停下来之后，家里安静了，但这个安静不是你想要的安静——这就是冷战。

\textbf{冷战的本质是：冲突从明码转成了暗码。}话停了，信道没关。

一个人不说话，并不等于没有输出。他吃饭的时间和平时不一样；他进门时看的是鞋柜不是你的脸；他把水杯放在桌子偏自己那一侧。\textbf{这些全都是发送。}区别只在于：明文的争吵你能逐句反驳，暗码的冷战你连反驳的对象都找不到——你不知道该回应哪一句，因为对方一句都没说。

真吵是双向的。冷战是\textbf{单方面把信道挂起}，同时让对方一直听见“挂起”这个动作本身。这是它比争吵更难熬的原因。

冷战里常见的四种载荷：

\begin{manstable}{P{1.6cm} L L}
\tbh{载荷} & \tbh{具体表现} & \tbh{传到对方那里是什么} \\
\midrule
时间 & 晚归、错开饭点、在卫生间待很久 & “我不想看见你” \\
空间 & 睡另一头、背对、把东西挪开 & “你被划出去了” \\
语言 & 只回“嗯”，不叫名字，用“那个”代替称呼 & “你已经不在我名单里” \\
动作 & 摔门、洗碗很响、把手机扣在桌上 & “我在表达，但我不给你抓手” \\
\bottomrule
\end{manstable}

需要说明的是，冷战看起来是省力的——你不用组织语言，不用面对冲突。\textbf{但它的成本比吵架高。}吵架是一次性消耗，冷战是持续供电。对方会一直在接收“我们没和好”这个信号，而你也没有真的在读别的什么。

还有一种更隐蔽的变体：\textbf{间歇性回复。}

\begin{mdquote}
他隔四个小时回你一句“在忙”。\par
过一天，又发来一条无关的新闻链接。
\end{mdquote}

这不是结束冷战，是维持冷战的一种更省电的方式。它比全沉默更黏，因为你每次都会重新开始等。需要说明的是，\textbf{这是冷战里最难判断的一种}：你分不清对方是在慢慢回温，还是在用最低的发送频率把你按在原地。

判据只有一条——\textbf{看对方有没有在回应正题。}回“在忙”、发无关链接，都是在回应你的等待，不是在回应那件事。\textbf{回应等待是维持，回应事情才是结束。}

那么冷战和冷处理的边界在哪？

\textbf{看它有没有长度。}冷处理是“我现在讲不清楚，给它两小时”。冷战是“我不知道要多久”。前者有终点，后者是把终点交给对方去猜。

\textbf{也看它有没有在回应。}冷处理期间，你仍然可以回应对方的\textbf{请求}（“我听到了，等我缓一下”），只是不处理那件事。冷战则连这条都不给。

\textbf{冷处理的上限通常是一到两天。}超过这个长度，它就不再是让情绪衰减，而是变成了惩罚——因为情绪本来就衰减不了那么久，多出来的时间是纯粹为了让人难受。

如果你发现自己正在冷战，唯一有效的动作还是 11.3 那一个：\textbf{让出发送权。}

\begin{mdquote}
“我现在还没想好怎么说，但我不想这样下去。给我一晚上。”
\end{mdquote}

这句话不认错，不复盘，只做一件事：\textbf{告诉对方这是暂停，不是终止。}

\section{吵完之后的三件事}

吵架停下来了。接下来有且只有三件事要做，而且顺序不能换：\textbf{道歉、复盘、翻篇。}

需要说明的是，\textbf{绝大多数关系不是被吵架拆掉的，是被这三件事的错序拆掉的。}先复盘再道歉，等于要求对方先认账；复盘之后不翻篇，等于把这件事永久留在桌面上，每次都能再拿出来用一次；只道歉不复盘，则等于约定下次照吵。

\textbf{第一件：道歉。}

道歉的作用不是认输，是\textbf{关闭这一局}。它要做的是让这件事从“未结”变成“已结”。一句合格的道歉有四个字段，缺一个就会失效。

\begin{manstable}{P{2.0cm} L L}
\tbh{字段} & \tbh{写什么} & \tbh{示例} \\
\midrule
指向 & 我做了什么 & 刚才打断你三次 \\
影响 & 它带来了什么 & 那让你没法把话说完 \\
不含条件 & 不能有“但是”“因为” & （后面什么都不接） \\
落点 & 下一步 & 下次你说的时候我先听完 \\
\bottomrule
\end{manstable}

\sample{1602}\par
\begin{mdquote}
\hciTag{反例}他说：「对不起，但我当时真的很累。」\par
——这是在道歉里塞了一个理由。对方的接收重点会落在“但是”上。他记住的不是你道歉了，是你还在解释。\par
\vspace{0.3\baselineskip}
✓ 他说：「刚才我打断了你三次。那让你没法把话说完。对不起。下次你先说。」
\end{mdquote}

值得注意的是，道歉有\textbf{时效}。冲突冷却后的\textbf{约十二到二十四小时}之间，是它最容易被接收的窗口。早于这个时候，你的道歉会被读成“想快点结束”；晚于这个时候，对方已经开始怀疑你是不是根本没打算说。

但也要防另一头：\textbf{不要用道歉当复盘的开场。}“我不是道歉了吗”这句话一出现，说明道歉已经被当成了一张入场券，而不是一次关闭。

\textbf{第二件：复盘。}

\begin{alertbox}[前提：只能在双方情绪落地之后进行。]
什么时候算落地？对方开始能谈论这件事的细节，而不再重复“你怎么能这样”。
\end{alertbox}

复盘不是把刚才那场又吵一遍。它要回答三个问题：

\begin{enumerate}
\item \textbf{这一局是怎么起来的？}不是“谁的错”，是“哪一句话让它升了级”。
\item \textbf{下次到哪一句就该停？}给一个具体的锚点，比如“说到‘你每次都这样’的时候停”。
\item \textbf{我们各自能改的一件事是什么？}只一件，而且要做得到。
\end{enumerate}

复盘有两个失败的典型。

第一个是\textbf{翻供}：

\sample{1603}\par
\begin{mdquote}
你：「那我们说好，下次一发现不对就先停一下。」\par
他：「那你得先承认刚才你也有问题。」\par
——复盘被拉回了追责。这一步一出现，前面所有谈好的内容全作废，两边又回到第一局。
\end{mdquote}

第二个是\textbf{把复盘做成判决}：全程在找谁的责任更大，最后得出一个结论，比如“这次你错七成，我三成”。\textbf{复盘要产出的是下一次的动作，不是这一次的分数。}

\textbf{第三件：翻篇。}

翻篇不是假装没发生过，是\textbf{约定这件事不再被调用}。

判断有没有真的翻篇，看一件事就够了：\textbf{下一次冲突里，这件事还会不会被拿出来当证据。}

\begin{itemize}
\item 如果会——那说明当时只是道歉了，没有翻篇；
\item 如果不会——才能说这一局真的关了。
\end{itemize}

翻篇需要一个明确的动作来标记，比如“这件事我放下了，不再提”，说出口。\textbf{不说出口的翻篇不算翻篇，只算暂时收起来。}

\section{常见场景：当众被顶}

需要说明的是，有一类冲突不是在两个人之间发生的，是在\textbf{有观众的时候}发生的。它有自己的一套机制。

\sample{1601}\par
\begin{mdquote}
周会，十二个人在座。汇报到一半，同事老周插话：\par
\vspace{0.3\baselineskip}
老周：「这个数据你上周不是说不是这个口径吗？」\par
你：「上周的口径是市场部后来改的，你看这封邮件——」\par
老周：「我不是跟你争这个，我是说大家别被带偏。」\par
你：「谁带偏了？」\par
\vspace{0.3\baselineskip}
会议室安静了两秒。这个你会开不下去了。
\end{mdquote}

拆开看：老周的第一句是内容之争（口径），第二句已经换成了状态之争（“你带偏大家”）——而且这个“大家”，把在场的十二个人全部拉进来当了证人。

\textbf{围观会把状态之争变成面子之争。}一旦有观众在场，任何一方退让，都不再只是让出发送权，而是当着十二个人的面让出。从此你要争的已经不光是“他有没有在听我”，是“这十二个人觉得我们谁是傻子”。

这时有几个动作是有效的：

\begin{itemize}
\item \textbf{停止当场辩。}在观众面前，你赢一句，成本是对方输一句的面子，这个账他要从别处讨回来。
\item \textbf{把议题拉回中性。}用“这封邮件我散会发给各位”这类动作，把争的东西从“谁对”换成“哪份文件”。
\item \textbf{不要用沉默回应沉默。}在当众场合里，你不说话，观众不会读成“你大度”，只会读成“你答不上来”。
\item \textbf{散会之后单独处理。}当众没解决的事，从来不会在当众解决。
\end{itemize}

\section{延伸：当对方是长辈、上级、陌生人或伴侣}

\textbf{当对方是长辈}：长辈的冲突里往往夹着\textbf{权力和历史}——他可能不是在对这件事发火，是在确认“你还听不听我的”。\textbf{此时先收的代价是双倍的}（你要让出发送权，还要让出“我是对的”这个位置）。可行的做法是\textbf{先收，再换个场合单独说事}。

\textbf{当对方是上级}：不要在当场争。上级在公开场合被顶，会把它处理成权威问题，而不是内容问题。\textbf{正确做法是先接住（“这个我回去想一下”），然后在下一次单独、私下的场合，用书面或数据把内容重新送一遍。}

\textbf{当对方是陌生人}：陌生人之间的冲突，绝大多数是内容之争。\textbf{能走就走，不要升级}——你们没有可供消耗的关系存量，一场冲突只能两败俱伤，没有任何东西值得赢。

\textbf{当对方是伴侣}：这是冲突里唯一\textbf{没有退出条款}的一类。

和陌生人不同，你不能走；和上级不同，你不能等下一次私下再说，因为你们睡在同一个房间；和长辈不同，权力关系是平的，谁也没有天然的裁决权。\textbf{代价是：普通人可以一走了之的冲突，在伴侣这里必须被处理。}

伴侣冲突里有两个反复出现的结构性故障。

\textbf{一是它会把当天的其他事拉进来。}今晚的争吵，通常会在二十分钟内连接到上周没洗的碗、上个月你没回的那条消息、去年过年你在他家的那次沉默。\textbf{原因不是他记性好，是当前这条载荷送不进去，系统在找别的载荷。}处理办法是先把当前这条送进去（先收），旧账会自己退场。

\textbf{二是它没有时间边界。}同事吵完可以各自回家，伴侣吵完还要一起吃饭、一起睡觉。这意味着冷处理的窗口被压缩得很短——\textbf{你走不掉，所以你必须比在其他关系里更早地做那个先收的人。}

\begin{mdquote}
常见的失效模式是\textbf{分房睡}。它看起来是给彼此空间，实际上是 11.6 里“空间载荷”的直接部署——它传过去的信息不是“我需要安静”，是“我从这张床上出去了”。\textbf{需要安静可以说出来，用不着搬走。}
\end{mdquote}

\section[常见误区]{常见误区}

\hcimistake{把“你冷静一下”当成降温}{\hciTag{反例}你：“你冷静一下，我们好好说。”}{这是对状态下达指令，无法执行。人在激动的时候听到“冷静”，只会更激动，因为这句话同时否定了他的情绪和你的接收动作。}

\hcimistake{以为把道理讲清楚就能结束}{\hciTag{反例}你：“你先听我把逻辑讲完，一二三四——”}{在状态之争里，每一个道理都在说“你说的不成立”，因此都在加重死锁。\hciSee{11.2}。}

\hcimistake{把旧账当成对方的攻击手段}{\hciTag{反例}他：“你去年也是这样的。”}{旧账被调出来的原因，是当前这条载荷送不进去。\textbf{把旧账压下去没用，要把当前这条送进去。}}

\hcimistake{以为沉默是在停火}{\hciTag{反例}你：「行，那我不说了。」（然后不回消息，不回房间，不接话）}{沉默也是发送。不说话的姿态、摔门的力度、背对的时长，全都是载荷。你以为你停了，其实你换了个信道在发。\hciSee{11.6}。}

\hcimistake{把“我先低头”理解成输}{\hciTag{反例}你：“行行行，都是我的错，行了吧。”}{先收不是认输，是先解锁。而“都是我的错，行了吧”不是先收，是换了一种姿势继续发送——它把认错做成了一件武器。}

\hcimistake{在有第三个人的场合继续}{\hciTag{反例}他（当着同事）：“这事你上周就答应了吧？”}{围观会把状态之争变成面子之争，从此更难收场。\hciSee{11.8}。}

\hcimistake{吵完之后立刻要一个结论}{\hciTag{反例}你：“那你以后到底改不改。”}{这个要求会把刚冷却的通道重新加热。它还要求对方在情绪未落地时，对一个还没有答案的问题表态。}

\hcimistake{把道歉和解释捆在一起发}{\hciTag{反例}你：“对不起，但我当时也是没办法。”}{“但是”之后的所有内容，都会被读成对前半句的撤回。对方记住的是“但是”，不是“对不起”。\hciSee{11.7}。}

%% ---------- 本章小结 ----------
\hciblocktitle{bullseye}{本章小结}

综上所述，本章的核心结论可以概括为以下几点：

\begin{enumerate}[label=\bfseries\arabic*., leftmargin=2.4em, labelsep=0.7em,
                  itemsep=0.45em, topsep=0.2em, parsep=0.2em]
\item 吵架是死锁：双方同时发送，没人接收。
\item 判断类型只看一句：争的是“怎么办”，还是“那句话什么意思”。八成是后者。
\item 状态之争里讲道理无效，因为道理解决的是内容。
\item 解锁只有一个办法：有一个人先收。代价是看起来像输，但死锁里没有赢家。
\item 冷战不是停火，是换了一条暗码信道接着发。它的成本比吵架更高。
\item 吵完之后有且只有三件事，顺序不能换：道歉、复盘、翻篇。
\item 冷战与冷处理的边界只有两条：有没有长度，有没有在回应正题。
\end{enumerate}

%% ---------- 练习 ----------
\begin{exercisessec}
\item 回忆你最近一次吵架（[请在此处填写当时的情境]）。逐句拆开，数一数有几句话是“接收”。

\item 在下一次争吵中，\textbf{对真人执行}：尝试在中途说一句“你先说”。不用道歉，不用让步，只说这一句。记录后续发生了什么。

\item 找出你生活里最近一次“错投”：你因为 A 事，对 B 发了火。写下真正该被路由的那一方。
\end{exercisessec}

%% ---------- 自测题 ----------
\begin{quizsec}
\item 对方在气头上，最不该说的一句话是？

\begin{enumerate}[label=\Alph*., leftmargin=2.2em, labelsep=0.6em,
                  itemsep=0.2em, topsep=0.3em, parsep=0.1em]
\item 你先说
\item 你冷静一下
\item 我听着
\item 你说完了我再讲
\end{enumerate}

\item 判断：争吵之后保持沉默，是让双方冷静下来的好办法。

\item 简答：为什么“先收”看起来像认输，实际上不是？
\end{quizsec}

%% ---------- 参考答案 ----------
\begin{answerssec}
\item B。它是对状态下达指令，无法执行，而且同时否定了对方的情绪和你的接收动作。\hciSee{11.10}。

\item 错。沉默也是发送。不说话的姿态、摔门的力度、背对的时长，全都是载荷。你以为你停了，其实你换了个信道在发；而这条信道更难判断，也更容易被读成惩罚。\hciSee{11.6}。

\item 因为死锁里没有“赢”这个结果——双方继续发下去，得到的只是两个人一起损失时间和关系。先收让出的只有发送权，没有承认任何指控。在死锁里，先收的人是掌握节奏的人，不是低头的人。\hciSee{11.3}。
\end{answerssec}

\hcirule

\begin{qabox}
\qaQ{读者提问}{如果对方就是想吵，根本不想解决，怎么办？}
\qaA{本手册回答}{这是一个非常好的问题。需要说明的是，有些人吵的不是这件事，而是这段关系本身是否要继续——那是另一个问题，本手册的这一章处理不了。你可以做的是：先把这次的载荷送进去（先收），如果仍然无效，那么问题已经不在通道上。}
\end{qabox}

\hcirule

\begin{qabox}
\qaQ{读者提问}{我和一个人已经冷战两周了。我该先开口吗？}
\qaA{本手册回答}{值得注意的是，冷战持续时间越长，开口的成本越高，但开口的收益不变。两周这个长度已经超过了冷处理的正常上限，也就是说，它现在的作用已经不是让情绪衰减，而是维持一个“我们还没和好”的状态。\textbf{先开口不等于先认错}——你只需要说“我不想这样下去了”，把暂停转成对话。至于那件事本身，那时再处理。搁置不是处理，只是把内容换了个地方存放。}
\end{qabox}

\hcirule

希望本章的内容能对你有所帮助。如需进一步了解第 12 章《时间》的相关内容，我可以随时为你展开。

%% 章末「页脚交互条」由 hci-manual.cls 自动挂接，无需手写。